% ---------------------------------------------------------------------------
\section{API Request and Response Example}
\label{app:api}
% ---------------------------------------------------------------------------

The request below was executed against the locally running API. The profile is
deliberately partial: the required fields are the neighbourhood, the overall quality,
the above-grade living area and the year built, and the optional fields supplied are
the remodelling year, overall condition, lot area, basement area, first-floor area,
bathrooms, bedrooms, rooms above grade and the garage.

{\footnotesize
\begin{verbatim}
curl -s -X POST http://127.0.0.1:8000/predict \
  -H "Content-Type: application/json" \
  -d '{"neighborhood": "CollgCr", "overall_qual": 6,
       "gr_liv_area": 1500, "year_built": 1975,
       "year_remod_add": 2005, "overall_cond": 5, "lot_area": 9000,
       "total_bsmt_sf": 800, "first_flr_sf": 900, "full_bath": 2,
       "half_bath": 0, "bedroom_abv_gr": 3, "tot_rms_abv_grd": 7,
       "garage_cars": 2, "garage_area": 480, "fireplaces": 0}'
\end{verbatim}
}

The response contains the point estimate, both conformal intervals, the confidence
classification and the assumptions applied to the omitted fields:

{\footnotesize
\begin{verbatim}
{
  "estimate": 146412.03,
  "currency": "USD",
  "interval_80": [130213.1, 164626.14],
  "interval_90": [124164.33, 172646.02],
  "confidence": "high",
  "confidence_score": 0.6742,
  "relative_interval_width_90": 0.3311,
  "model": "Elastic net",
  "model_version": "2026-11-19T13:20:58+00:00",
  "assumptions": {
    "sale_year": 2010,
    "sale_month": 6,
    "neighbourhood": "CollgCr",
    "derived_context_fields": ["amenity_density_1_5km",
      "amenity_score", "crime_index", "dist_hospital_km",
      "dist_park_km", "dist_school_km", "dist_transit_km",
      "neigh_lat", "neigh_lon", "park_area_ha",
      "school_rating", "walk_score"]
  },
  "generated_at": "2026-11-19T13:50:47+00:00"
}
\end{verbatim}
}

The arithmetic of the confidence class connects the statistical output to the
escalation policy. The 90 per cent interval half-width is
$(172{,}646.02 - 146{,}412.03)/146{,}412.03 = 0.179$, just inside the 0.18 threshold
for high confidence, and the score of 0.6742 follows the service's rule
$1 - 0.179/0.55$. The full relative width of 0.3311 is the quantity stored for
reporting and compared with the segment-level review thresholds recommended in
Section~\ref{sec:recommendations}.

{\footnotesize
\begin{verbatim}
[
  {"feature": "total_sf", "contribution_log": -0.0366,
   "direction": "decreases", "relative_effect_pct": 3.73},
  {"feature": "tot_rms_abv_grd", "contribution_log": -0.0351,
   "direction": "decreases", "relative_effect_pct": 3.57},
  {"feature": "bsmt_fin_type1", "contribution_log": -0.0252,
   "direction": "decreases", "relative_effect_pct": 2.55},
  {"feature": "quality_condition_index", "contribution_log": -0.0212,
   "direction": "decreases", "relative_effect_pct": 2.14}
]
\end{verbatim}
}

The explanation endpoint returns these local contributions for the same profile. All
four are negative because the property is modest relative to the comparison sample:
1,500 square feet of above-grade living area, a basement that the service defaults to
an unfinished finish type, and a quality--condition index of 30 place it below the
centre of its neighbourhood and vintage. The sign and magnitude are consistent with
the global drivers in Figure~\ref{fig:shap}, which is what makes a local explanation
useful to an agent: the attributes that drive the model globally also explain the
individual valuation, so the estimate can be discussed in the language of the listing
rather than the language of the model.
